\documentclass[11pt,a4paper,dvipdfmx]{article} 		% for platex and uplatex
\input{pieces/form00_header} % pieces
\input{pieces/kakenhi7} % pieces
\input{pieces/form01_header} % pieces
\input{pieces/form02_2023_header} % pieces
\input{pieces/hook3} % pieces
%#Name: ilr-2
\input{pieces/form04_jsps_eng_headers} % pieces
%\input{../../pieces/form04_mext_headers_2020} % pieces
\input{pieces/form04_ilr-2_header} % pieces
% ===== Global definitions for the Kakenhi form ======================
%　基本情報
%
%------ 研究課題名  -------------------------------------------
\newcommand{\研究課題名}{象の卵}

%----- 研究機関名と研究代表者の氏名-----------------------
\newcommand{\研究機関名}{逢坂大学}
\newcommand{\研究代表者氏名}{湯川秀樹}
\newcommand{\me}{\underline{\underline{H.~Yukawa}}} 
\input{pieces/inst_general} % pieces
\input{pieces/inst_ilr-2} % pieces
% user07_header
% ===== my favorite packages ====================================
% ここに、自分の使いたいパッケージを宣言して下さい。
\usepackage{wrapfig}
%\usepackage{amssymb}
%\usepackage{mb}
%\DeclareGraphicsRule{.tif}{png}{.png}{`convert #1 `dirname #1`/`basename #1 .tif`.png}
\usepackage{lineno}

% ===== my personal definitions ==================================
% ここに、自分のよく使う記号などを定義して下さい。
\newcommand{\klpionn}{K_L \to \pi^0 \nu \overline{\nu}}
\newcommand{\kppipnn}{K^+ \to \pi^+ \nu \overline{\nu}}

% ----- 業績リスト用 -------------
\newcommand{\paper}[6]{%
	% paper{title}{authors}{journal}{vol}{pages}{year}
	\item ``#1'', #2, #3 {\bf #4}, #5 (#6).			% お好みに合わせて変えてください。
}

\newcommand{\etal}{\textit{et al.\ }}
\newcommand{\ca}[1]{*#1}	% corresponding author;   \ca{\yukawa}  みたいにして使う
\newcommand{\invitedtalk}{招待講演}

\newcommand{\yukawa}{H.~Yukawa}					% no underline
%\newcommand{\yukawa}{\underline{\underline{H.~Yukawa}}}	% with 2 underlines
\newcommand{\tomonaga}{S.~Tomonaga}

\newcommand{\prl}{Phys.\ Rev.\ Lett.\ }		% よく使う雑誌も定義すると楽

% ===== 欄外メモ ==================
\newcommand{\memo}[1]{\marginpar{#1}}
%\renewcommand{\memo}[1]{}	% 全てのメモを表示させないようにするには、行頭の"%"を消す

%%\input{../../pieces/user07_header}	% expand
%\input{../../sample/simple/contents}	% skip
\input{pieces/hook5} % pieces

\begin{document}
\input{pieces/hook7} % pieces
%#Split: 01_framework  
%#PieceName: p01_framework
\input{pieces/p01_framework_00}
\section{Framework and Significance}
%　　　　＜＜最大　5ページ＞＞

%s02_framework_ilr
%begin frameworkILRWithInstructions ＝＝＝＝＝＝＝＝＝＝＝＝＝＝＝＝＝＝＝＝
\JSPSInstructions	% <-- 留意事項。これは消すか、コメントアウトしてください。

	The framework and significance of this international joint research
	are ...
%end frameworkILRWithInstructions ＝＝＝＝＝＝＝＝＝＝＝＝＝＝＝＝＝＝＝＝
\input{pieces/p01_framework_01}

%#Split: 02_plan  
%#PieceName: p02_plan
\input{pieces/p02_plan_00}
\section{Research Plan}
%　　　　＜＜最大　5ページ＞＞

%s08_plan_eng

%begin researchPlan ＝＝＝＝＝＝＝＝＝＝＝＝＝＝＝＝＝＝＝＝
	In the first year, we will visit major zoos in the world.
	We will search for the eggs in Africa in the second year, 
	and in India in the third year.

	More concrete plan is ...
%end researchPlan ＝＝＝＝＝＝＝＝＝＝＝＝＝＝＝＝＝＝＝＝

\input{pieces/p02_plan_01}

%#Split: 03_fostering  
%#PieceName: p03_fostering
\input{pieces/p03_fostering_00}
\section{Plan for Fostering}
%　　　　＜＜最大　2ページ＞＞

%s08_fostering_plan_eng

%begin fosteringPlan ＝＝＝＝＝＝＝＝＝＝＝＝＝＝＝＝＝＝＝＝
	Elephants incubate their eggs to raise their offspring.
	So should we.
	Any questions?
	
	\input{blahblah}
%end fosteringPlan ＝＝＝＝＝＝＝＝＝＝＝＝＝＝＝＝＝＝＝＝

\input{pieces/p03_fostering_01}

%#Split: 99_tail
\input{pieces/hook9} % pieces
\end{document}

